% from ega4-21c-fr.tex
\oldpage[IV-4]{323}
  \item Pour tout $x\in X$ qui n'est pas générisation d'un point de $S$ (autrement
  dit, tel que $\overline{\{x\}}\cap S=\varnothing$) et qui est tel que
  $\dim(\mathcal O_{X,x})\geq2$, l'anneau $\mathcal O_{X,x}$ est parafactoriel.
\end{enumerate}
\end{corollary}

\begin{proof}
Il est clair que l'ensemble $S'$ des points de $X$ générisations de points de $S$ est
l'intersection des voisinages ouverts de $S$. On peut se borner au cas où $X$ est
noethérien (les propriétés $a)$ et $b)$ étant locales sur $X$); comme tout cycle
$1$-codimensionnel sur $X$ qui est principal aux points de $S$ l'est aussi aux points
d'un voisinage ouvert de $S$, on voit que la condition $a)$ de (21.13.15) équivaut à
la condition $a)$ de (21.13.14) appliquée à la famille $(U_\lambda)$ des voisinages
ouverts de $S$. Le corollaire résulte alors de l'équivalence de $a)$ et $b)$ dans
(21.13.14).
\end{proof}

\begin{remark}[21.13.16]
\label{IV.21.13.16}
Sous les conditions générales de (21.13.14), supposons de plus que l'on ait
$\varinjlim\operatorname{Pic}(U_\lambda)=0$. Alors la condition $a)$ de (21.13.14)
est aussi équivalente à la suivante :
\begin{enumerate}
  \item[c)] \emph{Tout cycle $1$-codimensionnel sur $X$ dont le support est contenu
  dans un des ensembles $X-U_\lambda$ est localement principal.}
\end{enumerate}

On peut se borner au cas où $X$ est irréductible et tous les $U_\lambda$ sont non
vides. Il est clair que $a)$ entraîne $c)$, car si le cycle $1$-codimensionnel $Z$ sur
$X$ a son support dans $X-U_\lambda$, $Z|U_\lambda=0$, donc $Z|U_\lambda$ est
localement principal. Inversement $c)$ entraîne $a)$ : soit en effet $Z$ un cycle
$1$-codimensionnel sur $X$ tel que $Z|U_\lambda$ soit localement principal; comme $X$
est normal, $Z|U_\lambda=\operatorname{cyc}(D_\lambda)$, où $D_\lambda$ est un
diviseur sur $U_\lambda$ (21.6.10, (ii)), et l'hypothèse implique (en vertu de
(21.3.4)) qu'il y a un ensemble $U_\mu\subset U_\lambda$ tel que
$D_\mu=D_\lambda|U_\mu$ soit équivalent à $0$. Si
$D_\mu=\operatorname{div}(f_\mu)$, où $f_\mu$ est une fonction rationnelle régulière
sur $U_\mu$, $f_\mu$ peut être considérée comme une fonction rationnelle régulière
sur $X$. Si $Z'=\operatorname{cyc}(\operatorname{div}(f_\mu))$, on voit donc que
$Z''=Z-Z'$ a son support dans $X-U_\mu$; en vertu de l'hypothèse, $Z''$ est
localement principal, d'où la conclusion.

La condition $\varinjlim\operatorname{Pic}(U_\lambda)=0$ sera en particulier remplie
si $(U_\lambda)$ est la famille des voisinages ouverts d'un \emph{point} $z\in X$,
car pour tout $\mathcal O_{U_\lambda}$-Module inversible $\mathscr L_\lambda$, il
existe par définition un $U_\mu\subset U_\lambda$ tel que
$\mathscr L_\lambda|U_\mu$ soit isomorphe à $\mathcal O_{U_\mu}$. On voit donc que,
dans l'énoncé de (21.13.15), si $S=\{z\}$, les conditions $a)$ et $b)$ équivalent
aussi à la suivante :
\begin{enumerate}
  \item[c)] \emph{Tout cycle $1$-codimensionnel sur $X$ dont le support ne contient
  pas $z$ est localement principal.}
\end{enumerate}
Si de plus $\operatorname{Pic}(X)=0$, on en conclura que cette condition entraîne que
tout cycle $1$-codimensionnel dont le support ne contient pas $z$ est principal.
\end{remark}

\subsection{Le théorème de Ramanujam--Samuel}
\EGAReaderTarget{ega4.c04.subsec.21-14}{subsection:IV.21.14}
\label{IV.21.14}

\begin{theorem}[21.14.1]
\label{IV.21.14.1}
\textup{(Ramanujam--Samuel)}. ---
Soit $A$ un anneau local noethérien d'idéal maximal $\mathfrak m$, tel que son complété
$\widehat A$ soit intègre et intégralement clos. Soient $B$ un anneau local noethérien
tel que $\dim(B)\geq\dim(A)$, $\rho:A\to B$ un homomorphisme local faisant de $B$
une $A$-algèbre formellement lisse (pour les topologies préadiques
$(\mathbf 0, 19.3.1)$) et tel que le corps résiduel de $B$ soit fini sur celui de $A$.
Alors tout cycle $1$-codimensionnel sur $\operatorname{Spec}(B)$ qui est principal au
point $\mathfrak p=\mathfrak mB$, est un cycle $1$-codimensionnel principal.
\end{theorem}

\begin{proof}
Si $k=A/\mathfrak m$ est le corps résiduel de $A$, $B/\mathfrak mB$ est une
$k$-algèbre formellement lisse (pour sa topologie préadique) $(\mathbf 0, 19.3.5)$,
